\chapter{Project Management Plan}\label{chapter:manage}

\section{Team Work}

\subsection{Team Structure and Roles}
\begin{table}[H]
\centering
\small
\renewcommand{\arraystretch}{1.25}
\begin{tabularx}{\textwidth}{|p{3.2cm}|p{3.5cm}|X|}
\hline
\textbf{Role} & \textbf{Assignee} & \textbf{Core Responsibilities} \\
\hline
\textbf{Leader / AI Researcher} & Tran Dinh Gia \newline (SE181550) & 
$\bullet$ Lead the project and manage the schedule. \newline
$\bullet$ Formulate sequence modeling architectures and multi-task learning objectives. \newline
$\bullet$ Oversee model training and hyperparameter tuning. \\
\hline
\textbf{Edge Systems \& Model Optimization} & Truong Nhat Quang \newline (DE180544) & 
$\bullet$ Optimize decoding strategies (e.g., Beam Search). \newline
$\bullet$ Implement Knowledge Distillation and model quantization. \newline
$\bullet$ Benchmark inference latency and resource consumption on edge hardware. \\
\hline
\textbf{Data Engineering \& Backend Integration} & Nguyen Xuan Cuong \newline (DE180528) & 
$\bullet$ Construct and preprocess the large-scale Vietnamese parallel corpus. \newline
$\bullet$ Develop the data augmentation and synthetic noise generation pipelines. \newline
$\bullet$ Develop the offline inference backend and system integration APIs. \\
\hline
\end{tabularx}
\caption{Team Structure, Roles, and Responsibilities}
\label{tab:team_roles}
\end{table}

\subsection{Communication Plan}
\begin{itemize}
    \item \textbf{Daily Standups (Asynchronous):} Conducted via messaging platforms to track daily commitments and identify technical blockers.
    \item \textbf{Weekly Supervisor Consultation:} Formal weekly meetings with Ass.Prof Nguyen Luong Vuong to review experimental results, validate model performance, and plan subsequent phases.
    \item \textbf{Version Control:} Centralized GitHub repository for code collaboration, utilizing Pull Requests and peer reviews.
\end{itemize}

\section{Project Management Approach}

\subsection{Work Breakdown Structure}
\begin{table}[H]
\centering
\small
\renewcommand{\arraystretch}{1.2}
\begin{tabularx}{\textwidth}{|p{1.5cm}|p{3.8cm}|X|c|}
\hline
\textbf{Phase} & \textbf{Milestone} & \textbf{Key Deliverables} & \textbf{Timeline} \\
\hline
\textbf{Phase 1} & \textbf{Literature Review \& Data Pipeline} & Vietnamese text corpora collection, synthetic noise generation (diacritic stripping, space erasure), tokenization setup. & Weeks 1--4 \\
\hline
\textbf{Phase 2} & \textbf{Model Development \& Training} & Fine-tuning pretrained sequence-to-sequence model (T5-based), hyperparameter optimization, error analysis. & Weeks 5--9 \\
\hline
\textbf{Phase 3} & \textbf{Model Optimization \& Evaluation} & Decoding optimization, knowledge distillation, quantization, benchmarking BLEU, WER, and latency metrics. & Weeks 10--14 \\
\hline
\textbf{Phase 4} & \textbf{Backend \& System Integration} & Development of offline inference backend, API design for keyboard interaction, functional testing. & Weeks 15--17 \\
\hline
\textbf{Phase 5} & \textbf{Validation \& Deployment} & Comprehensive user testing, Keystroke Savings evaluation, final documentation, and project defense preparation. & Weeks 18--20 \\
\hline
\end{tabularx}
\caption{Work Breakdown Structure and Project Timeline}
\label{tab:wbs_table}
\end{table}

\subsection{Risk Management}
\begin{table}[H]
\centering
\footnotesize
\renewcommand{\arraystretch}{1.2}
\begin{tabularx}{\textwidth}{|p{2.5cm}|c|c|X|X|}
\hline
\textbf{Identified Risk} & \textbf{Sev.} & \textbf{Prob.} & \textbf{Impact} & \textbf{Mitigation Strategy} \\
\hline
\textbf{Domain Generalization Gap} & High & Med & High accuracy on training data but poor performance on noisy real-world typing. & Include synthetic typo noise (fat-finger errors) and non-standard words (teencode) during training. \\
\hline
\textbf{High Inference Latency} & High & Med & Model responds too slowly ($>50\text{ms}$), disrupting the typing experience. & Aggressively apply Knowledge Distillation and INT8 quantization; utilize lightweight Transformer variants. \\
\hline
\textbf{Compute Resource Limits} & Med & High & Slow training iterations due to large model sizes and extensive dataset. & Leverage distributed cloud GPUs and mixed-precision training (FP16). \\
\hline
\end{tabularx}
\caption{Risk Register and Mitigation Strategies}
\label{tab:risk_register}
\end{table}

\subsection{Quality Management}
\begin{itemize}
    \item \textbf{Linguistic Accuracy Metrics:} Evaluate using Word Error Rate (WER), Character Error Rate (CER), Exact Match ratio, and BLEU Score.
    \item \textbf{Performance Metrics:} Ensure inference latency remains under 50ms per keystroke event for seamless interaction.
    \item \textbf{Ergonomic Metrics:} Measure Keystroke Savings (KS) to quantify the reduction in physical typing effort compared to standard Telex.
\end{itemize}

\subsection{Budget and Funding (Optional)}
\begin{itemize}
    \item \textbf{Infrastructure:} Utilizing free-tier or academic allocations on platforms like Kaggle, Google Colab, and local workstations.
    \item \textbf{Total Budget:} \$0.00 USD.
\end{itemize}

\subsection{Change Management Process}
\begin{itemize}
    \item \textbf{Architecture Pivot:} If a chosen model architecture fails to meet latency requirements by Phase 3, the team will pivot to a smaller predefined baseline (e.g., CNN-LSTM hybrid).
    \item \textbf{Scope Adjustment:} Secondary features like personalized user dictionaries will be deprioritized if the core continuous typing engine requires extended development time.
\end{itemize}

\subsection{Closure and Evaluation}
\begin{itemize}
    \item \textbf{Deliverables Handover:} Source code, trained model weights, API documentation, and evaluation reports.
    \item \textbf{Final Thesis:} A comprehensive academic report detailing the methodology, empirical results, and system architecture.
\end{itemize}
